%% ============================================================
%%  chapter-template.tex
%%  Copy this into bible/chapters/ or spec/chapters/ as a starting
%%  point for a new chapter, then rename it and add its \include
%%  line to that document's main.tex.
%% ============================================================

\chapter{Chapter Title Here}

\section{First Section}

Body text goes here. Standard LaTeX — paragraphs, \emph{emphasis},
\textbf{bold}, footnotes\footnote{Like this.}, etc. all work normally.

\subsection{A Subsection}

Use \texttt{itemize} for bullet lists:

\begin{itemize}
  \item First point
  \item Second point
\end{itemize}

Use a note box to call out a definition, rule, or warning without
it competing visually with the body text:

\begin{sentinelnote}[Label Here]
Callout content goes here — a decision, a rationale, a rule.
\end{sentinelnote}

Use a table where useful:

\begin{table}[h]
  \centering
  \begin{tabular}{@{}ll@{}}
    \toprule
    \textbf{Column A} & \textbf{Column B} \\
    \midrule
    Value & Value \\
    \bottomrule
  \end{tabular}
  \caption{Table caption}
\end{table}

% Always end every chapter file with this — it renders the closing
% ornament and keeps a clean seam before the next \include'd chapter.
\sentinelchapterend
